\chapter{Correctness and Validation}

The performance results in Chapter 8 matter only if the system being measured is
doing the right work. For this thesis, correctness and measurement credibility
have several parts. Raft must order records correctly. Preferred leader election
must place leadership where Mako can submit work without weakening Raft safety.
Mako must receive committed records and replay them on followers. ReplayPool
must move replay work to parallel workers without changing the committed order.
Finally, the benchmark must report metrics that describe committed work,
follower progress, CPU use, and simulated-disk work consistently.

This chapter summarizes the validation evidence behind those claims. It is not
a test manual. Individual test names and helper details are less important than
the question each validation class answers.

\section{System and Fault Model}

The validation assumes the same three-replica structure used by the evaluation:
one submission replica and two backup replicas. Raft safety and availability
are based on a majority quorum, so two of the three replicas are enough to elect
a leader and commit new records. The experiments in this thesis colocate those
replicas on one machine and communicate over local TCP/IP. That choice keeps
the focus on Mako's replication backend, but it also means the validation is
not a wide-area network or datacenter-failure study.

The failure model is ordinary crash-stop and crash-recovery behavior for the
replication layer. A replica may be unavailable and later rejoin. The tests do
not assume Byzantine behavior, malicious replicas, disk corruption, or arbitrary
network partitions. The persistence work in Chapter 6 is also narrower than a
complete durable-recovery proof: it validates that modeled persistence work is
counted and reported, not that every possible storage failure has been tested.

One availability boundary is especially important for preferred leader
election. If the submission replica is unavailable, backup replicas can still
elect a Raft leader and preserve the replicated log according to Raft's normal
rules. That does not automatically mean the original Mako benchmark workload
continues unchanged, because the benchmark workers submit through the
leader-side submission role. Serving new transaction workload from a backup
leader would require the workload to be redirected or started there. This
thesis validates the Raft placement and recovery mechanism; it does not claim a
complete client redirection and production availability layer.

\section{What Validation Must Establish}

The first question is whether the Raft backend behaves like a consensus backend,
not just like a fast message path. It must elect leaders, replicate records,
commit only after enough replicas have accepted the record, and continue to make
progress across ordinary leader changes.

The second question is whether Raft fits Mako's execution model. Chapters 2 and
3 explained that Mako workers submit transaction log records through the
leader-side path, and Chapter 4 explained why Raft leadership must normally be
placed on that submission replica. Validation therefore needs to check both the
ordinary Raft behavior and the preferred-leader behavior added for Mako.

The third question is whether Single-Raft remains a real replicated design after
ReplayPool is added. ReplayPool should improve follower replay progress, but it
must not change the order Raft committed. It is valid only if committed records
are still delivered to the correct replay lane and replay preserves the ordering
required by that lane.

\section{Consensus and Leader Placement}

Standalone Raft validation covers the protocol-level behavior expected from the
backend: initial election, re-election, agreement on log entries, behavior
during failures, rejoin after a replica returns, and concurrent submissions.
Those tests do not prove the entire Mako system correct by themselves, but they
establish that the Raft layer is not merely a benchmark shortcut.

Preferred-leader validation covers the Mako-specific placement mechanism. The
tests and audits check that the submission replica is favored at startup, that a
backup replica can lead while the preferred replica is unavailable, and that
leadership is moved back only after the preferred replica has caught up. They
also check that preferred leadership remains compatible with log replication and
with the no-op and watermark behavior used to reason about committed progress.
The key point is the same as in Chapter 4: preference affects leader placement,
not Raft's voting or commit safety rules. The failback check is intentionally
guarded: the current leader attempts return to the preferred replica only after
observing that it has the committed prefix needed for safe leadership, and
Raft's ordinary voting rules remain the final guard.

\section{Mako Integration and Replay}

End-to-end Mako validation checks the path that matters for this thesis:
leader-side workers produce transaction log records, the replication backend
orders those records, and follower replicas replay the committed records into
their local state. This catches integration bugs that standalone Raft tests
cannot see, such as delivering a committed record to the wrong Mako lane or
measuring leader-side work while followers are not keeping up.

ReplayPool validation focuses on Single-Raft. The expected behavior is that Raft
decides the order, a lightweight delivery path hands committed records to Mako,
and follower replay workers perform the expensive database replay in parallel.
The validation checks are aimed at the boundary between those steps: records
from the same replay lane must remain ordered, records that cannot yet be safely
applied must wait, and follower replay counters must show whether replay is
actually progressing.

Chapter 5 also described the main caveat of parallel replay: records from
different replay lanes may finish in a different wall-clock order than their
global Raft log order. The validation therefore relies on Mako's replay safety
rules, including watermark checks and atomic state-update checks, to ensure that
parallel replay is safe for Mako's transaction records rather than assuming that
ReplayPool preserves one global replay completion order.

The practical invariant is that Raft still decides the committed order, while
Mako decides when each committed record is safe to apply to local state. Records
from the same lane are kept on the same replay path, so their lane order is
preserved. Records from different lanes may execute in parallel only under
Mako's existing replay checks. If a record is not yet safe for the local state,
it waits instead of being applied early.

\section{Persistence and Measurement}

Persistence validation is narrower than functional database recovery. It checks
that simulated-disk instrumentation is present and interpretable. Runs with
simulated persistence disabled should report disabled or zero storage work.
Runs with the NVMe-like or Cloud-SSD-like model enabled should report nonzero
persisted bytes and modeled writes. The sweep scripts must also carry those
values into the result CSV files so the graph data and the thesis text quote the
same measured work.

This validation supports the disk interpretation in Chapter 8. If a throughput
curve changes under simulated storage pressure, the thesis can point to the
corresponding bytes, writes, and bytes per committed transaction instead of
relying on throughput alone.

\section{Validation Matrix}

\tabref{tab:validation-matrix} summarizes the validation evidence used in the
thesis. The table intentionally groups tests by the claim they support rather
than listing every test helper.

\begingroup
\small
\setlength{\tabcolsep}{4pt}
\renewcommand{\arraystretch}{0.96}
\begin{longtable}{@{}>{\raggedright\arraybackslash}p{0.24\linewidth}
>{\raggedright\arraybackslash}p{0.46\linewidth}
>{\raggedright\arraybackslash}p{0.24\linewidth}@{}}
\caption[Validation evidence]{Validation evidence used to support the thesis claims.}
\label{tab:validation-matrix}\\
\toprule
Evidence class & What it checks & Thesis role \\
\midrule
\endfirsthead
\toprule
Evidence class & What it checks & Thesis role \\
\midrule
\endhead
Standalone Raft behavior & Leader election, re-election, agreement, failure cases, replica rejoin, and concurrent submissions. & Protocol sanity for the Raft backend. \\
\addlinespace
Preferred leader behavior & Startup placement, backup leadership while the preferred replica is absent, safe return after catch-up, log replication, and no-op progress. & Mako-compatible leader placement without changing Raft safety. \\
\addlinespace
Mako integration & Mako transaction logs are ordered, delivered, and replayed by follower replicas. & End-to-end replication correctness. \\
\addlinespace
ReplayPool behavior & Replay work moves to replay workers while same-lane ordering and deferred replay rules are preserved. & Single-Raft replay correctness and progress evidence. \\
\addlinespace
Persistence counters & Disabled-storage runs report no modeled work; simulated-storage runs report persisted bytes, modeled writes, and source breakdowns. & Disk-result interpretability. \\
\addlinespace
Benchmark harness & Leader throughput, worker CPU, follower replay progress, and simulated-disk counters are archived with the plotted results. & Measurement credibility. \\
\bottomrule
\end{longtable}
\endgroup

\section{Metric Discipline}

The evaluation uses leader-side committed throughput as the headline throughput
metric because it is the rate at which the benchmark completes committed
transactions on the submission side. That number is necessary, but it is not
sufficient. Single-Raft without enough replay parallelism can show high
leader-side throughput while follower replay falls behind, so follower replay
progress is reported as a separate sanity metric.

Worker CPU is interpreted similarly. High worker CPU helps show that Mako's
transaction workers are doing real work, but CPU utilization alone does not
prove replication correctness. Simulated-disk bytes and writes are also
separate from throughput: they explain how much modeled storage work the run
generated.

The plotted figures and the numbers quoted in the text are taken from the same
frozen result set. This matters because benchmark directories can be regenerated
as experiments are rerun. By freezing the result set used for the thesis, the
prose and graphs remain internally consistent: a throughput value discussed in
the text refers to the same measurement that produced the corresponding figure.
